% GENERATED FILE. Edit canonical lessons, then run npm run generate:books.
% canonical-source-hash: fnv1a64:78087ca84cb58d72
% canonical-lessons: SA-C422-prakarah, SA-C422-gopuram, SA-C422-parikha, SA-C422-suranga, SA-C422-khatam

\chapter{Walls and Gates}
\label{ch:sa-walls-and-gates}
\hlchaptermodality{422}

\begin{chapteropening}
\textbf{By the end of this chapter:} \emph{I can describe how an old town is walled and entered.}

Complete the last lesson of chapter 422: I can describe how an old town is walled and entered.
\end{chapteropening}

\section[\texorpdfstring{\sk{प्राकारः} (prākāraḥ)}{प्राकारः (prākāraḥ)}]{\sk{प्राकारः} (prākāraḥ) — a town wall; a rampart}

\label{lesson:SA-C422-prakarah}



\begin{quote}
Before the new one: say the Sanskrit for a mansion, then the Sanskrit for converses.
\end{quote}

\subsection*{You'll want to know: \sk{प्राकारः}}
\textbf{\sk{प्राकारः}} — \emph{prākāraḥ} — \textquotedblleft{}a town wall; a rampart\textquotedblright{}.

\begin{quote}
\sk{नगरं} \sk{परितः} \sk{उन्नतः} \sk{प्राकारः} \sk{अस्ति।}
\end{quote}

\begin{quote}
\emph{nagaraṃ paritaḥ unnataḥ prākāraḥ asti.}
\end{quote}

\begin{quote}
Around the city there is a high rampart.
\end{quote}

A masculine noun.

\begin{cousinweb}
\textbf{\sk{प्राकारः}} (\emph{prākāraḥ}) is built from \emph{prā}, a form of \emph{pra}, and the root √kṝ, \textquotedblleft{}to pour out, to heap\textquotedblright{}; MW gives \emph{prākāra} as a wall or enclosure, a rampart (MW).
\end{cousinweb}

\begin{grammarlens}[title={What you've built}]
One more word for getting around.
\end{grammarlens}

\subsection*{Your turn}
\begin{itemize}
\raggedright
  \item \emph{Say it:} \emph{prākāraḥ}
  \item \emph{Say it:} \emph{prākāraḥ}, once more
  \item \emph{Say it:} \emph{nagaraṃ paritaḥ unnataḥ prākāraḥ asti.} — Around the city there is a high rampart.
  \item \emph{Recall:} say the Sanskrit for a mansion, then the Sanskrit for converses, then say \emph{prākāraḥ} again
\end{itemize}

\begin{tcolorbox}[breakable,colback=teal!4,colframe=teal!35!black,title={Before you move on}]
What does \sk{प्राकारः} mean? (\textquotedblleft{}A town wall; a rampart\textquotedblright{}.) Say it once more.
\end{tcolorbox}
\section[\texorpdfstring{\sk{गोपुरम्} (gopuram)}{गोपुरम् (gopuram)}]{\sk{गोपुरम्} (gopuram) — a town gate}

\label{lesson:SA-C422-gopuram}



\begin{quote}
Before the new one: say the Sanskrit for converses, then the Sanskrit for a town wall.
\end{quote}

\subsection*{You'll want to know: \sk{गोपुरम्}}
\textbf{\sk{गोपुरम्}} — \emph{gopuram} — \textquotedblleft{}a town gate\textquotedblright{}.

\begin{quote}
\sk{यात्रिकाः} \sk{गोपुरेण} \sk{नगरं} \sk{प्रविशन्ति।}
\end{quote}

\begin{quote}
\emph{yātrikāḥ gopureṇa nagaraṃ praviśanti.}
\end{quote}

\begin{quote}
The pilgrims enter the city by the town gate.
\end{quote}

A neuter noun.

\begin{cousinweb}
\textbf{\sk{गोपुरम्}} (\emph{gopuram}) is built from \emph{go}, \textquotedblleft{}a cow\textquotedblright{}, and \emph{pura}, \textquotedblleft{}a town\textquotedblright{}; MW gives \emph{gopura} as a town gate, and in modern use it also names the tower over a temple gate (MW).
\end{cousinweb}

\begin{grammarlens}[title={What you've built}]
One more word for getting around.
\end{grammarlens}

\subsection*{Your turn}
\begin{itemize}
\raggedright
  \item \emph{Say it:} \emph{gopuram}
  \item \emph{Say it:} \emph{gopuram}, once more
  \item \emph{Say it:} \emph{yātrikāḥ gopureṇa nagaraṃ praviśanti.} — The pilgrims enter the city by the town gate.
  \item \emph{Recall:} say the Sanskrit for converses, then the Sanskrit for a town wall, then say \emph{gopuram} again
\end{itemize}

\begin{tcolorbox}[breakable,colback=teal!4,colframe=teal!35!black,title={Before you move on}]
What does \sk{गोपुरम्} mean? (\textquotedblleft{}A town gate\textquotedblright{}.) Say it once more.
\end{tcolorbox}
\section[\texorpdfstring{\sk{परिखा} (parikhā)}{परिखा (parikhā)}]{\sk{परिखा} (parikhā) — a moat}

\label{lesson:SA-C422-parikha}



\begin{quote}
Before the new one: say the Sanskrit for a town wall, then the Sanskrit for a town gate.
\end{quote}

\subsection*{You'll want to know: \sk{परिखा}}
\textbf{\sk{परिखा}} — \emph{parikhā} — \textquotedblleft{}a moat\textquotedblright{}.

\begin{quote}
\sk{प्राकारात्} \sk{बहिः} \sk{गभीरा} \sk{परिखा} \sk{अस्ति।}
\end{quote}

\begin{quote}
\emph{prākārāt bahiḥ gabhīrā parikhā asti.}
\end{quote}

\begin{quote}
Outside the rampart there is a deep moat.
\end{quote}

A feminine noun.

\begin{cousinweb}
\textbf{\sk{परिखा}} (\emph{parikhā}) is built from \emph{pari}, \textquotedblleft{}around\textquotedblright{}, and the root √khan, \textquotedblleft{}to dig\textquotedblright{}, which also gives \emph{khanati}, \textquotedblleft{}digs\textquotedblright{}; MW gives \emph{parikhā} as a moat or trench round a town or fort (MW).
\end{cousinweb}

\begin{grammarlens}[title={What you've built}]
One more word for getting around.
\end{grammarlens}

\subsection*{Your turn}
\begin{itemize}
\raggedright
  \item \emph{Say it:} \emph{parikhā}
  \item \emph{Say it:} \emph{parikhā}, once more
  \item \emph{Say it:} \emph{prākārāt bahiḥ gabhīrā parikhā asti.} — Outside the rampart there is a deep moat.
  \item \emph{Recall:} say the Sanskrit for a town wall, then the Sanskrit for a town gate, then say \emph{parikhā} again
\end{itemize}

\begin{tcolorbox}[breakable,colback=teal!4,colframe=teal!35!black,title={Before you move on}]
What does \sk{परिखा} mean? (\textquotedblleft{}A moat\textquotedblright{}.) Say it once more.
\end{tcolorbox}
\section[\texorpdfstring{\sk{सुरंगा} (suraṅgā)}{सुरंगा (suraṅgā)}]{\sk{सुरंगा} (suraṅgā) — a tunnel; an underground passage}

\label{lesson:SA-C422-suranga}



\begin{quote}
Before the new one: say the Sanskrit for a town gate, then the Sanskrit for a moat.
\end{quote}

\subsection*{You'll want to know: \sk{सुरंगा}}
\textbf{\sk{सुरंगा}} — \emph{suraṅgā} — \textquotedblleft{}a tunnel; an underground passage\textquotedblright{}.

\begin{quote}
\sk{दुर्गस्य} \sk{अधः} \sk{सुरंगा} \sk{अस्ति।}
\end{quote}

\begin{quote}
\emph{durgasya adhaḥ suraṅgā asti.}
\end{quote}

\begin{quote}
There is a tunnel beneath the fort.
\end{quote}

A feminine noun.

\begin{cousinweb}
\textbf{\sk{सुरंगा}} (\emph{suraṅgā}) is the word MW gives for a hole cut in a wall and an underground passage; MW cites the view that it is Greek \emph{syrinx}, \textquotedblleft{}a pipe\textquotedblright{}, so its origin is uncertain; it lives on as Hindi \emph{suraṅg}, a burglar's hole or a mine (MW; CDIAL 13506).
\end{cousinweb}

\begin{grammarlens}[title={What you've built}]
One more word for getting around.
\end{grammarlens}

\subsection*{Your turn}
\begin{itemize}
\raggedright
  \item \emph{Say it:} \emph{suraṅgā}
  \item \emph{Say it:} \emph{suraṅgā}, once more
  \item \emph{Say it:} \emph{durgasya adhaḥ suraṅgā asti.} — There is a tunnel beneath the fort.
  \item \emph{Recall:} say the Sanskrit for a town gate, then the Sanskrit for a moat, then say \emph{suraṅgā} again
\end{itemize}

\begin{tcolorbox}[breakable,colback=teal!4,colframe=teal!35!black,title={Before you move on}]
What does \sk{सुरंगा} mean? (\textquotedblleft{}A tunnel; an underground passage\textquotedblright{}.) Say it once more.
\end{tcolorbox}
\section[\texorpdfstring{\sk{खातम्} (khātam)}{खातम् (khātam)}]{\sk{खातम्} (khātam) — a ditch}

\label{lesson:SA-C422-khatam}



\begin{quote}
Before the new one: say the Sanskrit for a moat, then the Sanskrit for a tunnel.
\end{quote}

\subsection*{You'll want to know: \sk{खातम्}}
\textbf{\sk{खातम्}} — \emph{khātam} — \textquotedblleft{}a ditch\textquotedblright{}.

\begin{quote}
\sk{कृषकः} \sk{क्षेत्रस्य} \sk{समीपे} \sk{खातं} \sk{खनति।}
\end{quote}

\begin{quote}
\emph{kṛṣakaḥ kṣetrasya samīpe khātaṃ khanati.}
\end{quote}

\begin{quote}
The farmer digs a ditch near the field.
\end{quote}

A neuter noun.

\begin{cousinweb}
\textbf{\sk{खातम्}} (\emph{khātam}) comes from the root √khan, \textquotedblleft{}to dig\textquotedblright{}: it is the participle, \textquotedblleft{}dug\textquotedblright{}, and MW gives \emph{khāta} as dug, excavated and, as a noun, a ditch (MW).
\end{cousinweb}

\begin{grammarlens}[title={What you've built}]
One more word for getting around.
\end{grammarlens}

\subsection*{Your turn}
\begin{itemize}
\raggedright
  \item \emph{Say it:} \emph{khātam}
  \item \emph{Say it:} \emph{khātam}, once more
  \item \emph{Say it:} \emph{kṛṣakaḥ kṣetrasya samīpe khātaṃ khanati.} — The farmer digs a ditch near the field.
  \item \emph{Recall:} say the Sanskrit for a moat, then the Sanskrit for a tunnel, then say \emph{khātam} again
\end{itemize}

\begin{tcolorbox}[breakable,colback=teal!4,colframe=teal!35!black,title={Before you move on}]
What does \sk{खातम्} mean? (\textquotedblleft{}A ditch\textquotedblright{}.) Say it once more.
\end{tcolorbox}
